%% APPENDIX A -- spectral response, frequency conventions, primary-beam response.
%% Owner: appendix-triage.  Carries \label{app:conventions} (cited by
%% 04_method.tex) and \label{app:pbresp}.
\section{Spectral response, frequency conventions and the primary-beam
response}
\label{app:conventions}

%% The single-column figure of the three-point Hanning response against
%% sub-channel phase is removed; its numbers are in the peak-channel
%% paragraph below.  QUEUE_integrate.md entry S6: make_all.sh must stop
%% running make_fig_hanning.py and figorphan.py must carry
%% hanning_response.pdf in PENDING, or the figure gate fails.

\subsection{Frequency frames, linewidth and smearing}
%% (subsection label retired: nothing cites it)

Every frequency axis here -- searched spectra, tabulated ranges and
crossing frequencies alike -- is \emph{topocentric} as delivered by the
ALMA correlator, and the search applies no velocity-frame shift.
Velocities enter at the vetting stage, where the line mask carries every
crossing frequency to the barycentre at its own block's epoch and then to the
stellar frame before comparing it with a transition
(Appendix~\ref{app:linemask}). The residual error is then the catalogue
radial-velocity uncertainty alone, a few km\,s$^{-1}$ within 40\,pc;
second-order relativistic terms at the \MaskHalfKms\,km\,s$^{-1}$
scale are $\lesssim5$\,kHz at 350\,GHz, below the finest native channel
searched. The stacking is drift-aligned and inverse-variance weighted:
coherent in frequency trajectory, not in electric-field phase.

The native channel width is also the linewidth convention for every EIRP
threshold. The statistic is the single-channel flux density after drift
correction, so EIRP$_{\rm trig}$ is the minimum \emph{total} power of a
channel-confined carrier, flat in the assumed linewidth below the channel and
never a spectral power density. Above the channel it is not a threshold at
all: spreading a fixed total power over $N$ channels lowers the statistic as
$N^{\ScNBroadExp}$, measured on Eq.~\ref{eq:tstar} as the power recovered nine
times in ten, so a feature as wide as the widest line these fields resolve
(\ScNBroadChan{} channels) must carry \ScNBroadRatio$\times$ the power of a
carrier, where a filter matched to it would pay
\ScNBroadRatioMatched$\times$ and this search uses no such filter. The
Hz-resolution surveys' limits are total powers in the same sense, each the
best case of its own channelisation, which makes the axes of
Fig.~\ref{fig:context}(a) commensurable. Channel widths come from the
measurement sets' own CHAN\_WIDTH; MS~v2 defines no
``EFFECTIVE\_BANDWIDTH'' column, so the quoted channelisations are separations
and are not claimed to be noise bandwidths.

ALMA applies online Hanning smoothing by default \citep{ALMAHandbook}, so
the spectral response is ${\sim}2$ channels wide and neighbouring channels
are correlated. For the $0.25/0.5/0.25$ kernel the coefficients
are $\rho_1=\HanRhoOneFrac=\HanRhoOne$ and $\rho_2=\HanRhoTwo$, and the
smoothing narrows the noise by
$\sigma_{\rm smoothed}/\sigma_{\rm raw}=\HanSigRatio$. The peak
channel holds \HanBest{} of an unresolved carrier's flux at a channel centre
and \HanWorst{} on a boundary, a median \HanMed{} over sub-channel phase.
\emph{That loss enters the limits exactly once; the thresholds are free of it
and the limits are not.} A window's trigger flux is the bare
$\LtCommonSnr\sigma$ of its own spectrum, so the \emph{total} power of a
channel-confined carrier set against it would be optimistic by
$\times\HanFacBest$--$\times\HanFacWorst$. The injected carriers are not set
against it in that way: each is deposited with the correlator's own channel
response, at a sub-channel phase drawn afresh, and recovered through the same
trigger, which reads the peak channel. The measured completeness therefore
already contains the loss, and factorises into it exactly:
$\EirpNinetyMultA\,P_{\rm trig}$ is $\times\RpRespFac$ for the channel
response and $\times\RpMarginNinety$ for the margin a peak channel needs
above the nominal trigger for nine carriers in ten to return. So
$P_{\rm trig}$, the sensitivity $S_{\min}$ below and the filled symbols of
Fig.~\ref{fig:context}(a) are nominal throughout, which is what makes them
commensurable with other programmes' detection thresholds, while
$\mathrm{EIRP}_{90}$ carries the correction once.
The deposition falls \InjShortPct{} per cent short of the ideal response
averaged over phase, so the completeness is conservative by that much.

\emph{Intra-integration smearing, which the injections do not model.} Each
injected tone is deposited at one frequency per integration, so the campaign
carries the channel response and the drift between integrations but not the
sweep within one, where
$\eta_{\rm smear}$ is the mean of that response over the sweep and every
limit in this paper is divided by it once; the sweep cannot act inside the
transform, these channels being $\NmRespChanTau$ times wider than the
reciprocal of an integration, so it redistributes dwell time rather than
destroying coherence. The correction exceeds one per cent in
\NmRespNApplied, reaching $\times\SmcFacWorst$; it is measured where the
integration length is recorded, for \SmcNMeas{} of the \NWindows{} windows,
and \emph{bounded} for the other \SmcNBound{} at \SmcDtCeil\,s, the longest
integration delivered anywhere, where none falls below \SmcBoundEta.
Re-scoring the \SmcValNWin{} injected windows' own ladders with every tone
de-rated by its own $\eta_{\rm smear}$ gives a penalty never larger than the
one applied, so the correction errs the safe way.

\subsection{The primary-beam response}
\label{app:pbresp}

The sensitivity quoted for a window is $S_{\min}=5\sigma/A$, where $A$ is the
primary-beam response at the star's offset from the phase centre: the noise
$\sigma$ is measured in apparent flux density, the limit is quoted in true
flux density, and the two differ by $A$. Archival fields are
pointed at a catalogue position, often decades old, while these are
high-proper-motion stars, so the star can sit well off axis. Over the
\PbNWin{} windows the offset has a median of \PbOffMed$''$ and a
maximum of \BmOffMax$''$; the response $A$
runs from \BmRespMed{} in the median down to \BmRespMin, falling
below $0.9$ on \BmNBelowNine, every one of
those being a J2000 pointing on a fast-moving star rather than an error.

\emph{The response applied is a Gaussian of the full width the observatory
documents}, $\RingPbCoef\,\lambda/D$ \citep{ALMAHandbook}, evaluated at each
window's \emph{own} dish diameter. The dish matters: the \PbNAcaRow{} rows
taken on the \PbAcaRatio$\times$ smaller compact-array antennas have a
correspondingly wider beam, and a Gaussian forced to $D=12$\,m reproduces the
per-window geometry only to $\SensPbResidTwelve$ against $\SensPbResid$ for
the per-dish form. The blocked-aperture response, which carries the shadow of
the secondary and is what \textsc{casa} applies \citep{CASA2022}, has a
half-power width of $\BmAiryFwhm\,\lambda/D$ on a 12\,m antenna and
$\BmAiryFwhmAca\,\lambda/D$ on a 7\,m one; the Gaussian used here is that
beam to \BmDocAiryPct{} per cent in width, while a Gaussian of the Airy
first-null width departs from it in response by up to \BmModelMaxNullPct{} per
cent, which is why the documented width is the one adopted. The two forms part company as the
offset grows, which is what sets the retention rule of
\S\ref{sec:sample}, and \emph{the furthest
window of the primary census lies at $\BmCensusWorstC\,\lambda/D$, where they
agree to \BmCensusWorstPct{} per cent: that is the largest beam-model
uncertainty any sensitivity quoted here carries.}

\emph{One star is covered far off axis.} $\epsilon$\,Eridani, at
\BmEpsDistPc\,pc the nearest star the fine-channel search covers in
Band~6, has
\BmEpsNOn{} of its \BmEpsNWin{} windows at the pointing centre and
\BmEpsNOff{} lying $\BmEpsCLo$--$\BmEpsCHi\,\lambda/D$ off it, where the
response is \BmEpsRespLo--\BmEpsRespHi, inside the main lobe: the
blocked-aperture response does not reach its first null until
$\BmAiryNull\,\lambda/D$. \BmEpsNCensus{} of the \BmEpsNWin{} fall inside the
primary census, \BmEpsNCensusOff{} of them off axis, and carry its statistics,
counted in Table~\ref{tab:blockledger}; the deepest of them is the limit
Table~\ref{tab:persys} gives for this star; no window of it reaches the
trigger; and the \BmEpsNExt{} outside the census, \ExNOffAxis{} of them off
axis, are reported with the archival extension. Nor does any block of this star carry an annual-parallax record,
so its limit omits a displacement of \ExCenPxArcsec$''$ --
\ExCenPxBeams{} of that block's synthesised beam, from this star's parallax of
\BmEpsPlxMas\,mas, the largest of any the fine-channel search covers -- which
is smaller than the
\SensPxSmallBeams{} of a beam at which the retained amplitude below was
measured, so that measurement bounds the loss instead of being extrapolated
to it and the limit is optimistic by at most \BmEpsPxBoundPct{} per cent on
this account.

\emph{The correction is applied exactly once.} The offset is computed for
every measurement set from its own \texttt{FIELD::PHASE\_DIR} and the
propagated stellar position, and it enters in one place only:
$S_{\min}=5\sigma/A$ closes on all \PbSminClosed{} catalogued rows to
\PbSminWorstRes{} in relative terms, and on the \PbSminIndep{} of them whose
$A$ comes from the search's own stored geometry to \PbSminIndepRes{} -- the
other \PbSminTaut{} having $A$ back-solved from $5\sigma/S_{\min}$ and
flagged as such. Nothing else needs correcting: the statistic of
Eq.~\ref{eq:tstar}, the control ranks and the attribution are \emph{ratios}
formed inside one window, in which $A$ divides out identically.

\emph{The star's apparent direction includes the annual parallax, and the
amplitude a compact source retains when that is left out of the assumed phase
centre is measured rather than modelled.}
Re-extracting the \SensNReextWin{} windows of \SensNReextBlock{} blocks
furthest out at the full apparent position -- proper motion, parallax and
radial velocity where catalogued -- leaves each window's noise unchanged to
\SensReextNoiseMaxPct{} per cent and alters no disposition. A carrier
injected at the apparent position and recovered from the \emph{same}
visibilities at both does return the normalised dirty beam, with no model in
it: \SensPxSmallMeas{} against a predicted \SensPxSmallPred{} at
\SensPxSmallBeams{} of a beam, and \SensPxLargeMeas{} against
\SensPxLargePred{} at \SensPxLargeBeams. The measurement exceeds the model at
the larger displacement, so a limit corrected by the model is shallower than
it need be.

\emph{A component of a bound pair does not travel in a straight line, and the
propagation assumes it does.} Of the \OmNStar{} searched positions,
\OmNComponent{} have a companion inside \OmRadius{} arcsec
\citep{wds,orb6} and \OmNIndOnly{} more have a disturbed astrometric solution
with no companion found. Where the pair has a published visual orbit the
displacement follows from it, where it has not the projected separation puts a
ceiling on the orbital acceleration with the companion's mass taken from
\citet{pecaut2013}, and the Hipparcos--Gaia proper-motion anomaly
\citep{kervella2022} covers \OmNPma{} of them. The largest
displacement anywhere is \OmWorstArcsec{} arcsec at \OmWorstStar,
\OmWorstBeams{} of that block's synthesised beam, no crossing lies more than
\OmCrossWorstBeams{} of a beam from its predicted position and the deepest
limit of any system moves by at most \OmSysMovedMaxPct{} per cent -- except
for \OmFlagStars, whose solutions are disturbed with no companion located and
no Hipparcos epoch to measure one from, so that the \OmNWinFlag{} windows of
those \OmNFlag{} stars carry limits optimistic by an unmeasured amount.

%% Carried forward: the UV Ceti worked example (flux -> EIRP) is out of this
%% appendix because two of its four numbers had no macro behind them.  If a
%% worked example is wanted back it needs three: the flux density, the implied
%% EIRP and the distance used.  It is the clearest single verification aid the
%% paper could carry.  \label{app:repro} and the transmitter-power benchmark
%% are both resolved: nothing cites the label any more, and exactly one
%% benchmark survives, in Sec. 6.
%% tab:blockledger lives here rather than in Sec. 3: it is a full-width
%% table* worth 0.53 pp of the main text, and Sec. 3 carries its headline
%% rows in prose and cites it.
\begin{table*}
\centering
\scriptsize
%% ★ v4.26 (r25-captions): 4 printed lines -> the author's 3.  DELETED because
%% \S\ref{sec:sample}'s first paragraph says it in the same words -- "every
%% sensitivity, every expectation and every limit here is computed over it"
%% -- checked before deleting.  No macro and no \ref lost a citation site.
\caption{\textbf{The survey ledger, part one: execution blocks and windows.}
Each row carries its own number, and a group sums to the row whose number
heads it. Hours are window-hours, so the classes partition the hours as they
partition the windows. Rows~3.2 to 3.5 are the archival gap, over which no
result is formed.}
\label{tab:blockledger}
%% GENERATED by blockledger_v411.py -- do not hand-edit.
\begin{tabular}{@{}r@{~}p{0.33\textwidth}rrrrr@{}}
\hline
 & Set of execution blocks & Blocks & Windows & Stars & Systems & Hours \\
\hline
1 & Archival scope: blocks the query exposes & 656 & -- & -- & -- & -- \\
1.1 & \hspace{0.8em}\ldots{} processed with the same pipeline & 479 & -- & -- & -- & -- \\
1.2 & \hspace{0.8em}\ldots{} not processed at the freeze & 177 & -- & -- & -- & -- \\
\hline
2 & Processed with the same pipeline$^{c}$ & 484 & -- & -- & -- & -- \\
2.1 & \hspace{0.8em}\textbf{Primary census} & 407 & 1663 & 90 & 83 & 1098 \\
2.1.1 & \hspace{1.6em}Class~A, fine channels$^{a}$ & -- & 405 & 66 & 61 & 228 \\
2.1.2 & \hspace{1.6em}Class~B, coarse channels$^{a}$ & -- & 1258 & 89 & 82 & 870 \\
2.2 & \hspace{0.8em}Reserved hold-out & 77 & 315 & 36 & -- & 214 \\
\hline
3 & Not processed at the freeze: what became of each & 177 & -- & -- & -- & -- \\
3.1 & \hspace{0.8em}Calibrated and searched & 144 & -- & -- & -- & -- \\
3.1.1 & \hspace{1.6em}\textbf{Archival extension}, reported & 115 & 463 & 16 & 16 & 269 \\
3.1.1.1 & \hspace{2.4em}Class~A, fine channels$^{a}$ & -- & 67 & 11 & -- & 48 \\
3.1.2 & \hspace{1.6em}Searched after the extension was frozen & 8 & 32 & -- & -- & -- \\
3.1.3 & \hspace{1.6em}$\epsilon$\,Eri, both pointings$^{d}$ & 8 & 32 & -- & -- & -- \\
3.1.4 & \hspace{1.6em}No window reached the export & 13 & 0 & -- & -- & -- \\
3.2 & \hspace{0.8em}No calibration in the ALMA delivery & 18 & -- & -- & -- & -- \\
3.3 & \hspace{0.8em}Calibration failed & 11 & -- & -- & -- & -- \\
3.4 & \hspace{0.8em}Not attempted & 3 & -- & -- & -- & -- \\
3.5 & \hspace{0.8em}Working disk insufficient & 1 & -- & -- & -- & -- \\
\hline
4 & Beyond 40\,pc: out-of-sample control & 47 & 198 & 22 & -- & 88 \\
5 & Rank-calibration sample$^{b}$ & 326 & 1322 & 56 & -- & -- \\
\hline
\end{tabular}

{\footnotesize\raggedright $^{a}$The two classes partition the census windows, not its blocks: 255 blocks carry a Class~A window and 392 a Class~B one, and most carry both, so the hours of the two rows partition the windows and not the observing time.\\ $^{b}$Not an out-of-sample set and not an addition to this ledger: every one of its blocks is counted in a row above, in the primary census or in the hold-out, and it is listed apart only to keep it from being counted twice.\\ $^{c}$Panels 1 and 2 overlap rather than add: 479 of these 484 blocks are row 1.1, and the other 5 carry no member observing unit in the archive metadata, so the scope query does not expose them; 4 of the 5 are in the primary census.\\ $^{d}$Every one is an $\epsilon$\,Eri Band~6 block. The archival pointings put the star either at the centre of the field, at a response of 1.00, or well off it, where it is 0.23--0.28 over the 16 of this row's 32 windows that lie off axis, above the 0.22 retention floor of \S\ref{sec:sample} in every case. All of them are searched, each from its own recorded phase centre, and their conventions are in Appendix~\ref{app:conventions}.\\ The same band's 12 windows that fall in the primary census are counted in row~2.1 and carry every census statistic; the on-axis pointings reach a response of 1.00.}

\end{table*}

Table~\ref{tab:blockledger} accounts for every execution block the archival
query exposed and Table~\ref{tab:counts} for every event the search returned;
the full itemisation of both, all \MtCountRows{} rows, accompanies this paper
as data, and the row numbers below are that itemisation's. The telescope
on-source time behind the primary census is
\StTelHours\,h, of which \StTelHoursFine\,h carries at least one Class~A
window and \StTelHoursCoarse\,h at least one Class~B window -- figures that
overlap and must not be added. Confirmed planet counts there and in the
target list come from the \SfHostCatalogue{} \SfHostTable{} table
\citep{NEA2013}, accessed \SfHostDate. Three counts in the second table are
\MtCrossA{} over three different sets -- the Class~A crossings of the adopted
ledger, those of the per-window extraction catalogue and the quality-passing
crossings -- and the first two differ on \FrNFlag{} rows, where the statistic
read over a narrower field falls on the other side of the trigger from the
whole-primary-beam reading the ledger adopts, in both directions
(Appendix~\ref{app:falsealarm}); the accompanying table carries the ledger's
flag, \MtCrossACen{} of them in the census. Row 3.1.1 applies the criterion
of \S\ref{sec:summary} and row 2.2.1 is decomposed system by system in
Appendix~\ref{app:falsealarm}.

\begin{table}
\centering
\scriptsize
\caption{\textbf{The survey ledger, part two: events, dispositions and
confirmation.} Every headline count in this paper is read from here.}
\label{tab:counts}
%% GENERATED by counts_r15.py -- do not hand-edit.
%% The accounting rows of the master count table; the itemisation is in
%% the machine-readable file this generator writes beside it.
\begin{tabular}{@{}r@{~}p{0.405\columnwidth}r@{\,\,}p{0.325\columnwidth}@{}}
\hline
 & Count & Value & What it counts \\
\hline
1 & Threshold crossings at $T_\star\ge5$ & \EvNCrossRaw & census and archival tail \\
1.1 & Class~A, fine channels & \MtCrossA & the carrier experiment \\
1.2 & Class~B, coarse channels & \LdgNCrossB & all in the one block failing \S\ref{sec:bdblock} \\
\hline
2 & Quality-passing crossings & \EvNCrossQp & the same rows as 1.1 \\
2.1 & \ldots{} attributed to a transition & \EvNAttrQp & of \MkNCoinc{} coincidences inside the mask in the star's own frame \\
2.2 & \ldots{} unattributed & \EvNUnattrQp & rows 2.2.1 to 2.2.3 \\
2.2.1 & \ldots{} in the primary census & \CeObs & against \CeExp{} expected by chance at the \CeNWin{} windows' own control positions -- a count chance delivers with probability \CeProbNull \\
2.2.2 & \ldots{} in archival-tail blocks & \CeNTailCross & outside that comparison \\
2.2.3 & \ldots{} coincident with an atmospheric transition & \CeTlObsCut & within \ScRadMHz\,MHz of one in the topocentric frame (Appendix~\ref{app:rfi}), and counted as atmospheric on both sides of that comparison \\
\hline
3 & Class~A systems & \NSysClassA & of \NSystems{} surveyed \\
3.1.1 & Class~A systems with same-frequency coverage in a separated session & \textbf{\ZcConfSysDay} & \textbf{the only confirmation count quoted}: a later block in another session covering a frequency the first covered, \ZcConfSysYr{} of them over a year later. For an accelerating carrier, confirmation further requires the repeat to fall at a matching orbital phase. Over the census together with the repeat blocks the recurrence test reads \\
3.3 & \ldots{} with no separated second block at a covered frequency & \ZcNoRepSys & confirmation impossible over the \ZcNRecEB{} blocks the recurrence test reads; \RcNoRepSysAll{} of all \NSystems \\
\hline
4 & Class~A union bandwidth & \CvUnionA\,GHz & \DomAIslands{} disjoint islands \\
4.1 & \ldots{} removed by the line mask & \MkMaskAPct\,\% & \MkMaskAGHz\,GHz: the figure the text quotes, against \FrMaskPct{} per cent of the \DnuAB\,GHz archival union \\
\hline
\end{tabular}

\end{table}

%% ★★ 2026-10-08: THE ALMA PROJECT CODES, RESTORED AS A TABLE.
%% They were taken out of the Acknowledgements as padding and then lived "in
%% the data release"; there is no data release, so they were nowhere.  ALMA's
%% data-use policy expects the codes to be traceable from the publication, so
%% this is not optional -- but an appendix table satisfies the policy without
%% putting \NProjCodes{} codes back into the front matter, which is what was
%% objected to.  The fragment is generated by projcodes_v400.py and is the
%% same list that resolves every block in Table~\ref{tab:blockledger}.
\begin{table}
\centering
\scriptsize
%% ★ v4.26 (r25-captions): 4 printed lines -> the author's 3.  The unbreakable
%% \mbox is about a fifth of a column here, so this caption breaks badly and
%% costs a whole line more than its character count suggests -- which is why
%% the limit is measured in printed lines and not in characters.
\caption{\textbf{The \NProjCodes{} ALMA proposals the searched blocks belong
to}, as \mbox{ADS/JAO.ALMA\#\textit{code}}. Every census block belongs to
one.\ProjCodeUnresClause}
\label{tab:projcodes}
%% GENERATED by projcodes_v400.py -- do not hand-edit.
\mbox{2012.1.00142.S}, \mbox{2012.1.00198.S}, \mbox{2012.1.00268.S}, \mbox{2012.1.00385.S}, \mbox{2012.1.00437.S}, \mbox{2012.1.00993.S}, \mbox{2013.1.00359.S}, \mbox{2013.1.00588.S}, \mbox{2013.1.00645.S}, \mbox{2013.1.01147.S}, \mbox{2013.1.01214.S}, \mbox{2015.1.00307.S}, \mbox{2015.1.00633.S}, \mbox{2015.1.00783.S}, \mbox{2015.1.01015.S}, \mbox{2015.1.01260.S}, \mbox{2016.1.00104.S}, \mbox{2016.1.00358.S}, \mbox{2016.1.00515.S}, \mbox{2016.1.00637.S}, \mbox{2016.1.00803.S}, \mbox{2016.1.00817.S}, \mbox{2016.1.00878.S}, \mbox{2016.1.00880.S}, \mbox{2016.2.00015.S}, \mbox{2016.2.00200.S}, \mbox{2016.A.00013.S}, \mbox{2017.1.00167.S}, \mbox{2017.1.00200.S}, \mbox{2017.1.00215.S}, \mbox{2017.1.00561.S}, \mbox{2017.1.00698.S}, \mbox{2017.1.00704.S}, \mbox{2017.1.00786.S}, \mbox{2017.1.00986.S}, \mbox{2017.1.01583.S}, \mbox{2017.1.01644.S}, \mbox{2018.1.00045.S}, \mbox{2018.1.00072.S}, \mbox{2018.1.00126.S}, \mbox{2018.1.00461.S}, \mbox{2018.1.00470.S}, \mbox{2018.1.00947.S}, \mbox{2018.1.01149.S}, \mbox{2018.1.01545.S}, \mbox{2018.1.01833.S}, \mbox{2019.1.00189.S}, \mbox{2019.1.00696.S}, \mbox{2019.1.01443.T}, \mbox{2019.1.01517.S}, \mbox{2019.1.01603.S}, \mbox{2019.1.01681.S}, \mbox{2019.2.00208.S}, \mbox{2021.1.00485.S}, \mbox{2021.1.01209.S}, \mbox{2021.2.00097.S}, \mbox{2022.1.00134.S}, \mbox{2022.1.00338.L}, \mbox{2022.1.00793.S}, \mbox{2022.1.01163.S}, \mbox{2023.1.00390.S}, \mbox{2023.1.00487.S}, \mbox{2024.1.00165.S}, \mbox{2024.1.00595.S}, \mbox{2024.1.00722.S}, \mbox{2024.1.01095.S}, \mbox{2024.A.00040.S}

\end{table}
